% Chapter 14: The Doorway
% Scene function: Post-turn horror set piece. Visual manifestation
% (mirror technique). Gordon responds by opening the sickroom door
% for the first time. The haunting pushes Gordon toward grief, not
% away from faith. The sickroom thread resolves through Gordon.

\chapter{The Doorway}

Week four. Direct manifestation.

This is the phase I have been building toward since the first night in the basement. Everything before this was preparation: the ambient wrongness, the sub-threshold erosion of Gordon's relationship with the house's predictability. All of it was foundation. Week four is when the foundation becomes visible.

The protocol calls for a visual element. Not a full apparition. Atmospheric Hauntings uses visibility the way a painter uses shadow: not as the subject, but as the thing that gives the subject depth. A shape at the edge of perception. A figure in peripheral vision that is gone when you turn. A reflection in a surface that shows something the room does not contain.

The technique is positioning. You have to know exactly where the subject's eye will be, and when, and for how long. The bathroom mirror, for instance: if I stand in the hallway at a specific distance from the door frame while Gordon is brushing his teeth, there is a quarter-second window when his eye, tracking his own reflection, will pass through the angle where the hallway is visible behind him. A quarter second. If I am in position, he will see something. If I am not, he will not. The margin is zero.

On Thursday night I set up for the mirror.

Gordon goes to the bathroom at 11:08. I am in the hallway, six feet back from the door frame, at the angle I have calculated from three weeks of observation. The bathroom light is on. The hallway is dark. When Gordon looks at the mirror, the light differential will make the hallway behind him briefly visible, and I will be in it.

He brushes his teeth. He looks at the mirror.

I do not know what he sees. I know what I am presenting: a shape, upright, at the far edge of what the mirror contains. Not a face. Not a figure. A suggestion of something that occupies space in a way that empty hallways do not. The eye registers it before the mind can name it, and by the time the mind looks, I am no longer there.

Gordon stops brushing. He looks at the mirror for a long moment. Then he turns and looks down the hallway.

The hallway is empty. I have moved.

He stands in the bathroom doorway looking at the dark hall for about fifteen seconds. This is a long time to look at an empty hallway. Then he does something I did not expect.

He walks to the sickroom door.

He has not opened this door since I have been in this house. He has walked past it every night without pausing. I have watched him walk past it dozens of times, and the not-pausing was so consistent that I stopped noting it in the shift reports. It was a fixed behavior. Part of the baseline.

He opens the door.

He stands in the doorway. He does not go in. He looks at the room where his wife died. The room I have not used. The room I said I was holding in reserve.

The curtains are still sheer. The depression in the floor is still there. The nightstand with the evaporated water glass is still there. Nothing has changed since I opened this door myself, weeks ago, and cataloged it as an environmental feature.

Gordon is not cataloging it. He is standing in the doorway of the room where he took care of his wife, and the light from the hallway falls across the depression in the floor where the bed was.

He stands there for a long time. Then he closes the door, gently, and goes to bed.
